\hypertarget{appendix-f1a6d43e80d1}{}\label{referencestart:appendix-f1a6d43e80d1}
\section{Screen-space ray generation and mouse picking}\label{appendix-f1a6d43e80d1-screen-space-ray-generation-and-mouse-picking}

\hypertarget{appendix-f1a6d43e80d1-screen-space-ray-generation-and-mouse-picking}{}

\textbf{Curriculum depth:} Developed · \textbf{First seen:} Phase 11 · \textbf{Latest development:} Phase 11

\begin{allmightroledefinition}

\subsection{Definition and mechanism}\label{appendix-f1a6d43e80d1-definition-and-mechanism}

\hypertarget{appendix-f1a6d43e80d1-definition-and-mechanism}{}

A screen ray begins with a pixel location mapped through the viewport to NDC, then through inverse projection and camera pose into a world-space origin and direction. For perspective cameras, rays usually share the eye origin; orthographic rays have parallel directions but different origins across the image. Mouse picking must use the same high-DPI scaling, viewport offset, origin, and Y convention as rendering. Near and far unprojection can define a ray segment, but a point at infinity or an infinite far plane needs an explicit direction construction. Picking then intersects that ray with scene geometry or uses an object-ID buffer.

\end{allmightroledefinition}

\subsection{Prerequisites}\label{appendix-f1a6d43e80d1-prerequisites}

\hypertarget{appendix-f1a6d43e80d1-prerequisites}{}

\begin{itemize}
\tightlist
\item
  \href{../04-Light-Transport-and-Image-Synthesis/main.pdf\#nameddest=appendix-1c07d64f5e70}{Camera view and clip-space pipeline (Light Transport and Image Synthesis, Reference notes: Camera view and clip-space pipeline)}
\item
  \hyperref[appendix-f744e41b95b4]{Unit directions and safe normalization}
\item
  \hyperref[appendix-7d356c396d1a]{Visual angle and calibrated rays}
\end{itemize}

\subsection{Used by}\label{appendix-f1a6d43e80d1-used-by}

\hypertarget{appendix-f1a6d43e80d1-used-by}{}

\begin{itemize}
\tightlist
\item
  No downstream concept node is added in this stage.
\end{itemize}

\subsection{Full teaching explanation}\label{appendix-f1a6d43e80d1-full-teaching-explanation}

\hypertarget{appendix-f1a6d43e80d1-full-teaching-explanation}{}

\begin{itemize}
\tightlist
\item
  \href{../03-Geometry-Cameras-and-Image-Formation/main.pdf\#nameddest=02-foundations-ii-ch035-recover-geometry-from-depth-and-construct-screen-space-rays}{Recover geometry from depth and construct screen-space rays (Geometry, Cameras, and Image Formation, Ch. 4)}
\end{itemize}

\subsection{Source coverage}\label{appendix-f1a6d43e80d1-source-coverage}

\hypertarget{appendix-f1a6d43e80d1-source-coverage}{}

Supplied lesson parts: 1082--1083.

\subsection{Future keywords}\label{appendix-f1a6d43e80d1-future-keywords}

\hypertarget{appendix-f1a6d43e80d1-future-keywords}{}

inverse projection, viewport mapping, picking ray
